\documentclass[11pt,a4paper]{article}

\usepackage[utf8]{inputenc}
\usepackage[T1]{fontenc}
\usepackage{lmodern}
\usepackage[margin=1in]{geometry}
\usepackage{booktabs}
\usepackage{enumitem}
\usepackage{xcolor}
\usepackage{hyperref}
\usepackage{fancyhdr}
\usepackage{array}

\hypersetup{
  colorlinks=true,
  linkcolor=blue!60!black,
  urlcolor=blue!60!black,
  pdftitle={Rubik's Cube Solver --- User Manual},
  pdfauthor={Akash Chakraborty and Atal Chaudhuri},
}

\newcommand{\key}[1]{\fbox{\texttt{\strut#1}}}
\newcommand{\app}{\textsc{RubiksCubeSolver}}

\pagestyle{fancy}
\fancyhf{}
\fancyhead[L]{\app{} --- User Manual}
\fancyhead[R]{\thepage}
\renewcommand{\headrulewidth}{0.4pt}

\title{\vspace{-1.5em}\Huge \app\\[0.3em]\Large User Manual}
\author{Akash Chakraborty \and Atal Chaudhuri\\[0.5em]
\small Sister Nivedita University, Kolkata}
\date{\today}

\begin{document}
\maketitle
\thispagestyle{empty}

\begin{center}
\fbox{\parbox{0.85\textwidth}{\centering
\textbf{About this manual.} This guide walks you through installing and
using the \app{} application that ships as a single-file executable on
Windows and a \texttt{.app} bundle on macOS. The application implements
the algorithms described in \emph{Multi-Solution Search and the Limits
of Neural Guidance in Rubik's Cube Solving}; this manual covers only the
user-facing workflow, not the algorithmic details.
}}
\end{center}

\vspace{1em}
\tableofcontents
\vspace{1em}

%===========================================================================
\section{What This Application Does}
%===========================================================================

\app{} is an interactive desktop tool that accepts a 3-by-3 Rubik's Cube
state from the user and returns a move sequence that solves it. Under the
hood it runs an enhanced version of Kociemba's Two-Phase Algorithm with
four selectable solver modes:

\begin{itemize}[leftmargin=2em]
  \item \textbf{Baseline (K=1)} --- the original Kociemba solver, fastest.
  \item \textbf{Multi-Solution (K=5)} --- our primary contribution; explores five
        Phase~1 endpoints and keeps the shortest combined solution. On
        average cuts 0.53 moves off the baseline and roughly doubles the
        fraction of $\leq$20-move solutions.
  \item \textbf{Anytime (K=5, 2s)} --- the multi-solution solver with a 2-second
        soft budget, returning the best solution found within the budget.
        Designed for interactive applications where latency matters.
  \item \textbf{Neural LUT (K=5)} --- multi-solution search augmented with a
        precomputed lookup-table move ordering. Included for
        completeness; quality matches plain K=5.
\end{itemize}

Input is either painted onto a 3D cube visualisation, loaded from one of
four preset demo scrambles, or generated as a random scramble. Output is
a human-readable move list that can also be played back one step at a
time on the 3D view.

%===========================================================================
\section{Getting the Application}
%===========================================================================

\subsection{Windows}

Download \texttt{RubiksCubeSolver.exe} (single file, roughly 50~MB).
Double-click to launch. No Python or other dependencies are required;
everything is bundled inside the executable.

\textbf{First launch.} The first time you run the executable, Windows
extracts bundled assets into a temporary directory. This takes a few
seconds. Subsequent launches are faster.

\textbf{SmartScreen warning.} Because the executable is unsigned, Windows
Defender SmartScreen may show a blue warning dialog on first launch.
Click \textit{More info} $\to$ \textit{Run anyway}.

\subsection{macOS}

Download \texttt{RubiksCubeSolver.app.zip}, double-click to unzip, and
drag the resulting \texttt{RubiksCubeSolver.app} to your
\texttt{Applications} folder.

\textbf{First launch (Gatekeeper).} Because the app is unsigned, double-
clicking it the first time will show \emph{``cannot be opened because
Apple cannot check it for malicious software''}. Work around this by
\emph{right-clicking} the app and choosing \textit{Open}; macOS will
then show an Open button you can click through. Subsequent launches work
from a normal double-click.

%===========================================================================
\section{The Main Window}
%===========================================================================

When you launch \app{} you see four panels laid out in a single window:

\begin{itemize}[leftmargin=2em]
  \item \textbf{3D view (top-left).} A rotatable 3D rendering of the cube.
        Left-drag to rotate, right-click a sticker to select it.
  \item \textbf{2D unfolded net (top-middle).} The standard cross layout
        showing all six faces at once. The currently selected sticker is
        highlighted.
  \item \textbf{Info panel (top-right).} Displays the selected sticker,
        the current solver mode, the control reference, and (after a
        solve) the move count and solution summary.
  \item \textbf{Colour palette (bottom-left).} Six swatches: White, Red,
        Blue, Yellow, Orange, Green. The currently selected colour is
        highlighted.
\end{itemize}

Four yellow \textit{Demo} buttons and a \textit{Random Scramble} button
sit near the bottom-right. A single-line status message at the bottom of
the window shows the most recent action or result.

%===========================================================================
\section{Providing the Cube State}
%===========================================================================

There are four ways to tell the solver what cube to solve. Pick whichever
is most convenient.

\subsection{Demo scrambles (recommended for first-time users)}

Click any of the four yellow \textit{Demo} buttons. Each loads a
preselected valid cube state drawn from the paper's seed-42 benchmark.
The 3D view and unfolded net update immediately and the status bar
reports \textit{``Random scramble applied. Press Enter to solve!''}
These are convenient for comparing the different solver modes on
identical input: load the same demo twice, solve once with each mode,
and compare the move counts.

\subsection{Random scramble}

Click \textit{Random Scramble} (or press \key{S}) to apply 20--30 random
moves to a solved cube. The result is a valid, uniformly distributed
scramble. Use this when you want a fresh, unpredictable cube.

\subsection{Painting stickers manually}

If you want to solve a specific physical cube in your hand, you can
paint the state one sticker at a time:

\begin{enumerate}[leftmargin=2em]
  \item Orient your physical cube with \textbf{blue on top, red facing you, yellow on the right}
        (the orientation shown in the app's header text).
  \item Click a sticker on the 3D view or unfolded net to select it. The
        selected sticker is highlighted.
  \item Pick a colour, either by clicking a swatch in the palette or by
        pressing a digit \key{1}--\key{6} (1:W, 2:R, 3:B, 4:Y, 5:O, 6:G).
        The selected sticker repaints instantly.
  \item Press \key{Tab} / \key{Shift+Tab} to advance to the next /
        previous face, or use the arrow keys to move within a face.
\end{enumerate}

Centre stickers are locked to prevent a common input error (they define
which colour belongs to which face and must match the physical cube's
orientation).

\subsection{Guided mode}

Press \key{G} to toggle \textit{guided mode}, which walks you through
the cube sticker by sticker in a fixed order. Useful the first time you
enter a cube by hand --- each press of a number key paints the current
sticker and advances automatically. Press \key{G} again to leave guided
mode.

%===========================================================================
\section{Choosing a Solver Mode}
%===========================================================================

The info panel shows the currently active solver mode below the cursor
position. Press \key{M} to cycle through the four modes:

\begin{center}
\begin{tabular}{@{}ll p{0.50\textwidth}@{}}
\toprule
\textbf{Mode} & \textbf{Label} & \textbf{What it does} \\
\midrule
1 & Baseline (K=1)        & Kociemba's original solver. Very fast
                            ($\sim$10~ms on this build), slightly longer solutions. \\
2 & Multi-Solution (K=5)  & Our primary algorithm: tries 5 Phase 1
                            endpoints, returns the shortest. Default mode. \\
3 & Anytime (K=5, 2s)     & Multi-solution search that always returns within
                            2~seconds using the best solution found so far. \\
4 & Neural LUT (K=5)      & Multi-solution search with a lookup-table move
                            predictor. Quality matches K=5. \\
\bottomrule
\end{tabular}
\end{center}

The \textbf{default mode is Multi-Solution (K=5)}. If you are just trying
to solve a cube quickly and do not care about the algorithmic
comparison, leave this alone.

%===========================================================================
\section{Solving the Cube}
%===========================================================================

Once a valid cube state is in place, press \key{Enter}. The solver runs
on the current mode. Depending on the mode, a solve takes between
$\sim$10~ms (baseline) and $\sim$15~s (K=5 on a hard cube). During a
solve the window is unresponsive; wait for the status bar to update.

When the solve completes:

\begin{itemize}[leftmargin=2em]
  \item The info panel shows the mode used, the move count, and the
        solution string.
  \item The status bar at the bottom reports success or failure.
  \item If no solution could be found within the internal time cap
        (15~s for K=5, 2~s for Anytime), you will see
        \textit{``Error: Timeout, no solution within given time.''}
        Re-running the solve or switching to Baseline almost always works.
\end{itemize}

\subsection{Stepping through the solution}

Press \key{N} to advance one move. The 3D view animates the rotation and
the status bar describes the move in plain English (for example
\textit{``U prime: top face counter-clockwise''}). Keep pressing \key{N}
until the cube is solved. This is the core feature for using the app as
a physical solver: make the move on your cube, then press \key{N} for
the next one.

\subsection{Comparing two modes on the same cube}

A common demonstration workflow:

\begin{enumerate}[leftmargin=2em]
  \item Click \textit{Demo 1} to load a fixed scramble.
  \item Press \key{M} until the info panel reads \textit{Baseline (K=1)}.
  \item Press \key{Enter}. Note the move count (e.g. 21).
  \item Click \textit{Demo 1} again to reset the same cube.
  \item Press \key{M} until the info panel reads \textit{Multi-Solution (K=5)}.
  \item Press \key{Enter}. Note the move count (e.g. 18).
\end{enumerate}

The 3-move saving on Demo~1 and Demo~2, or 2-move saving on Demo~3 and
Demo~4, is the visible, reproducible demonstration of our contribution
on a single cube. Across 200 random cubes the mean saving is 0.53
moves ($p < 10^{-20}$, paired test).

%===========================================================================
\section{Keyboard and Mouse Reference}
%===========================================================================

\subsection{Solving and mode}
\begin{center}
\begin{tabular}{@{}ll@{}}
\toprule
\textbf{Key} & \textbf{Action} \\
\midrule
\key{Enter}  & Solve the current cube with the active mode. \\
\key{N}      & Step to the next move in the solution. \\
\key{M}      & Cycle through the four solver modes. \\
\key{R}      & Reset the cube to solved. \\
\key{S}      & Apply a random scramble (equivalent to the button). \\
\key{Q}      & Quit the application. \\
\bottomrule
\end{tabular}
\end{center}

\subsection{Input and navigation}
\begin{center}
\begin{tabular}{@{}ll@{}}
\toprule
\textbf{Key} & \textbf{Action} \\
\midrule
\key{Tab}          & Move to the next face. \\
\key{Shift+Tab}    & Move to the previous face. \\
\key{$\uparrow$} \key{$\downarrow$} \key{$\leftarrow$} \key{$\rightarrow$}
                   & Move the cursor within the current face. \\
\key{1}--\key{6}   & Paint the selected sticker with the corresponding colour. \\
\key{G}            & Toggle guided-mode input. \\
\bottomrule
\end{tabular}
\end{center}

\subsection{Mouse}
\begin{itemize}[leftmargin=2em]
  \item \textbf{Left-drag on the 3D view:} rotate the view.
  \item \textbf{Click on a sticker (3D or 2D):} select that sticker.
  \item \textbf{Click on a colour swatch:} set the selected colour.
  \item \textbf{Click a Demo or Random Scramble button:} load a scramble.
\end{itemize}

%===========================================================================
\section{Troubleshooting}
%===========================================================================

\paragraph{``Windows protected your PC'' dialog on first launch.}
The executable is unsigned. Click \textit{More info} $\to$ \textit{Run
anyway}. Subsequent launches do not show this.

\paragraph{``Cannot be opened because Apple cannot check it'' on macOS.}
Right-click the \texttt{.app} and choose \textit{Open}. Then click
\textit{Open} in the follow-up dialog. This is a one-time click-through.

\paragraph{The app takes five seconds to show its window on Windows.}
Expected. PyInstaller single-file executables extract their contents to
a temporary folder on first launch, which takes a few seconds. The
window appears as soon as extraction finishes.

\paragraph{``Error: Timeout, no solution within given time.''}
The current cube is hard enough that the solver did not complete within
the mode's time cap. Options:
\begin{itemize}[leftmargin=2em]
  \item Press \key{M} to switch to Baseline and press \key{Enter} again
        --- Baseline uses the C backend and solves in milliseconds.
  \item Click the same Demo button again to reload the cube, then
        retry; solve time varies slightly run-to-run.
\end{itemize}

\paragraph{``Cannot change center piece!''}
Centre stickers are locked because they define which colour is assigned
to which face. If the physical cube you are entering has a different
centre colour arrangement than what the app shows, re-orient the physical
cube to match (blue on top, red facing you, yellow on the right).

\paragraph{The 3D view and the unfolded net look out of sync.}
They are always derived from the same internal state. If they look
different, the most likely explanation is that the 3D view has been
rotated --- only the camera changed, not the cube. Left-drag to rotate
it back.

\paragraph{I pressed Enter and nothing happens.}
Check the status bar. If it shows an error about colour counts, the
painted cube is invalid (there must be exactly nine stickers of each
colour). Use \key{R} to reset and try again, or click a Demo button to
load a known-valid state.

%===========================================================================
\section{Where the Algorithm Comes From}
%===========================================================================

The solver logic and the four modes correspond directly to sections of
the accompanying paper, \emph{Multi-Solution Search and the Limits of
Neural Guidance in Rubik's Cube Solving}. Readers who want the
algorithmic detail behind a given mode should consult:

\begin{center}
\begin{tabular}{@{}ll@{}}
\toprule
\textbf{Mode in app} & \textbf{Paper section} \\
\midrule
Baseline (K=1)        & Section 2.2 (Two-Phase Algorithm) \\
Multi-Solution (K=5)  & Section 4.3 and Algorithm 1 \\
Anytime (K=5, 2s)     & Section 4.4 \\
Neural LUT (K=5)      & Section 4.5 and Section 6.8 \\
\bottomrule
\end{tabular}
\end{center}

This user manual does not duplicate that material. If you just want to
solve a cube, everything you need is in Sections 1--7 above.

\end{document}
